\section{Provenance-Aware Data and Evidence Pipeline}
\label{sec:evidence}

The data pipeline turns independently collected payloads into frozen experimental units while preserving the distinction between signal values, collection states, provenance, and verified experiment facts. This separation is central to the paper's claims: the system must be able to reconstruct why a record entered an experiment without exposing the same metadata to the component being evaluated.

\subsection{Release and Schema Freezing}

Every formal run records the collector release, App schema, Browser schema, Web-probe revision, backend revision, relation catalogs, and the hashes of the corresponding configuration files. The active paired-data path uses FeatureApp version code 8, a fixed 177-signal App contract, a fixed 67-signal Browser contract, and a versioned pair-provenance schema. The current controlled App-side attack bundles were produced with a later debug-oriented FeatureApp release that preserves the \appfp signal contract while adding intervention controls needed by the runner. The paper treats these releases as distinct cohorts rather than pretending that the current App attack pairs and the development \pairedfp snapshot were collected under one identical binary.

This distinction does not prevent the present \appfp analysis because the evaluated fields and their states follow the frozen 177-signal contract. It does, however, prevent the existing data from supporting a joint paired attack claim. Before the future external-browser evaluation, the collection and intervention controls must be unified in one release, and both the baseline and active phases must attempt Browser capture under the same App, Browser, and backend schemas.

\subsection{Snapshot Construction}

Snapshot construction begins from canonical raw payloads rather than from convenience projections. For each candidate App record, the builder verifies the release and schema, recomputes or checks the canonical payload hash, follows the collection receipt, confirms the associated batch state, and compares any available analysis projection with the authoritative raw data. For a candidate Browser pair, it also verifies the Browser schema and payload, both sides of the provenance link, the pair-event sequence, and the completed terminal status. Selection decisions are deterministic and are recorded rather than being applied silently.

A successful run produces the signal views together with the artifacts needed to audit them. These artifacts include the frozen feature catalog, an App/Browser sample index, a per-input selection audit, quarantine records with reason codes, quality-control summaries, and a dataset manifest that binds the run to its inputs and configuration. The current development/QC snapshot contains 26 accepted App observations. Seventeen have completed Browser provenance and enter \pairedfp; nine are retained as App-only because the Browser side did not complete. The snapshot is unlabeled and exists to validate collection, pairing, missingness, and runtime-input construction. It is not used as evidence of attack-detection effectiveness.

The App-only path is not a degraded 244-dimensional record. It is an explicit 177-signal data product with a control-plane explanation of why no paired Browser record exists. Likewise, quarantine is not a catch-all deletion path: it preserves which contractual condition failed, allowing a collector, backend, or pairing defect to be distinguished from a legitimate but incomplete observation.

\subsection{Evidence Extraction}

For each accepted App observation, the extractor first verifies that the value and state maps match the frozen catalog. It then parses version strings, platform families, graphics descriptions, display geometry, memory and resource values, sensor and battery context, and host-topology fields into normalized facts. Parsing is conservative. A relation premise is considered available only when the source field has an appropriate state and the parser can produce the semantic component required by the relation.

Cross-layer comparisons are derived after this normalization step. The extractor may, for example, compare the Android host family with the runtime User-Agent and platform, compare a WebView provider major version with the runtime Chromium major version, examine whether the FeatureApp bridge topology is present, or classify Native and Web graphics strings into compatible platform families. Each comparison retains the underlying field paths and any tolerance or applicability context. Relations involving coarse or unstable values are not forced into exact equality; they remain soft, contextual, or unassessed when the available evidence does not justify a stronger result.

The extractor also removes control-plane metadata before constructing the runtime input. Provider names, experiment phases, attack tools, labels, stable identities, receipt identifiers, and split assignments never become evidence fields. The final EvidenceBundle is serialized deterministically and hashed, so the same payload, field-state map, and extractor version produce the same evidence representation. This property supports caching, run comparison, and later verification of every cited fact.

Field states allow the pipeline to distinguish records that would look identical after conventional imputation. An observed API value of zero is evidence about the runtime, whereas a timeout says only that the runtime did not return a value under the collection conditions. Similarly, an Android-version relation is not scored as an extreme mismatch when the active User-Agent no longer contains an Android token. The extractor records that the premise became unavailable or inapplicable, and the relation evaluator must respect that state.

\subsection{Relation Lifecycle}

A plausible natural-language association does not become an executable relation automatically. Each candidate must specify its premise fields, parsing logic, applicability condition, comparison direction, severity, tolerance, counterexamples, and expected behavior on normal observations. It must also identify its primary knowledge source. Official documentation can establish that a field represents an Android release, a WebView provider, a User-Agent surface, or a display metric; the project must still justify the cross-layer inference that connects those meanings.

The current official-derived catalog therefore contains both executable relations and deliberately deferred candidates. Coarse semantic contradictions that do not require learned thresholds can be compiled, including Android host versus desktop/script surface incompatibility, compatible Android-version chains, WebView-provider/runtime major-version relations, default/runtime User-Agent compatibility, expected FeatureApp bridge topology, development-configuration context, and explicit Android-versus-Windows graphics contradictions. By contrast, display geometry, memory, sensor counts, battery behavior, provider policy, and future attestation evidence require larger normal datasets, temporal measurements, device-family baselines, deployment policy, or fields not yet collected. They remain candidates rather than being assigned arbitrary thresholds.

Device-mined predicates follow a parallel but separate lifecycle. Their evidence records must identify the training or exploratory data from which they were derived, the independent data used to challenge them, and known normal counterexamples. A predicate is compiled only after its field paths and execution semantics have been reviewed. The two catalogs can express similar conclusions, but preserving source identity makes it possible to determine whether coverage arises from documented semantics, empirical regularity, or both.

\subsection{Data Leakage Controls}

Verified experiment facts are joined to signal records through control-plane identifiers and canonical hashes. This join determines whether a record is eligible for a task and which baseline-active pair it belongs to, but the joined labels and tool metadata are not inserted into the signal payload or EvidenceBundle. The evaluator processes the baseline and active payloads independently; configuration and mutation annotations are reattached only after relation states have been frozen for reporting.

The split protocol also operates on stable groups and scenarios rather than individual rows. Observations that may share a physical device, provider device profile, controlled scenario, or paired manipulation remain in the same partition. This prevents repeated sessions or nearly identical attack templates from appearing on both sides of an evaluation boundary. Group-aware weighting and evidence-hash deduplication can reduce the influence of repeated profiles during training, but they do not replace group-separated evaluation.

These controls address several practical shortcut channels. A model can otherwise learn provider-specific strings, experiment directory names, runner-injected identifiers, tool names embedded in raw payloads, or fold labels rather than the intended environment evidence. By freezing a signal-only representation and retaining provenance in separate manifests, HybridGuard can audit the complete experiment without giving the decision path access to the answer.
